\chapter{Introduction}
\label{chap:introduction}

\section{Background and Motivation}
\label{sec:intro_background}

Corporate social responsibility (CSR) is often discussed as a commitment by firms to account for the social and environmental consequences of their operations. In supply chains, however, CSR is not implemented by statements alone. It has to be translated into concrete decisions: who invests, who coordinates, who bears cost, who receives the economic benefit, and how independent firms sustain cooperation when their private incentives are not naturally aligned. This implementation problem is especially important in supply chains because the firm that makes or benefits from a CSR promise is not always the firm that performs the operational work required to make the promise credible.

This dissertation studies CSR as a strategic coordination problem in supply chains. Its central premise is that CSR becomes economically meaningful only when it is embedded in pricing decisions, investment incentives, channel power, reverse-logistics design, and cost-sharing rules. A manufacturer may benefit from a stronger responsible brand position, but a retailer may be closer to consumers and may bear the practical burden of take-back service. A system-wide recycling network may reduce environmental harm, but its facilities, transportation flows, and operating costs have to be allocated across stakeholders. These tensions make CSR a question of governance and coordination, not only a question of ethics or reputation.

EV battery recycling provides the case context for this dissertation because it makes the operational content of CSR unusually visible. As electric vehicle adoption grows, end-of-life lithium-ion batteries create both an environmental risk and a resource-recovery opportunity \citep{Harper2019, HoarauLorang2022}. Responsible battery management requires collection, sorting, traceability, reverse transportation, processing, and material recovery. It is also shaped by regulatory pressure such as extended producer responsibility (EPR), which places responsibility for end-of-life management on producers while the physical return process often depends on retailers, service centers, and downstream recycling infrastructure \citep{yang2025bridging, Gu2021, ZhangW2024}. In this setting, CSR is not an abstract preference. It appears as recycling maturity, consumer trust in take-back systems, compliance credibility, and fair responsibility for shared reverse-channel costs.

For this reason, EV battery recycling is treated here as an applied case study for a broader CSR coordination problem. The dissertation does not ask only how batteries should be recycled. It asks how CSR commitments can be organized so that independent supply-chain actors have economic incentives to invest in responsible recycling, coordinate across organizational boundaries, and share the cost of common infrastructure.

\section{Research Problem}
\label{sec:intro_problem}

The core research problem is that CSR commitments in supply chains often create benefits and costs at different points in the channel. A CSR-oriented recycling program may increase consumer demand, strengthen brand credibility, and help satisfy regulatory expectations, but the investment may be made by a different actor than the one that captures the largest share of the benefit. In a bilateral manufacturer--retailer channel, this creates a vertical incentive problem. In a multi-retailer reverse-logistics system, it creates a network-design and cost-allocation problem.

The EV battery recycling case highlights both problems. First, a manufacturer and a retailer may jointly benefit from a more credible battery take-back and recycling system, but they may disagree about who should lead the channel and who should pay for recycling maturity. Second, once many retailers collect batteries across a region, CSR responsibility becomes a shared infrastructure problem: the system must decide where batteries flow, which facilities are used, and how the resulting cost is divided. A centralized physical optimum is not enough if independent stakeholders regard the allocation as unfair or unstable.

The overarching research question is therefore:

\begin{quote}
\emph{How can CSR commitments in supply chains be translated into economically viable coordination mechanisms, using EV battery recycling as a case study?}
\end{quote}

This question connects the two studies developed in Chapters~\ref{chap:model1} and~\ref{chap:model2}. Study~1 focuses on CSR-driven vertical coordination between one manufacturer and one retailer. Study~2 extends the CSR responsibility problem to a multi-retailer reverse-logistics network and examines how shared recycling costs can be allocated fairly.

\section{Research Objective and Questions}
\label{sec:intro_objectives}

The objective of this dissertation is:

\begin{quote}
\emph{To develop a game-theoretic framework for analyzing how CSR responsibilities can be coordinated, financed, and sustained in supply chains, with EV battery recycling serving as the empirical and operational case context.}
\end{quote}

The thesis is guided by one overarching question and two study-specific questions.

\begin{itemize}
    \item \emph{How can CSR commitments be operationalized in supply chains so that independent firms have incentives to invest, coordinate, and share responsibility?}

    \item \emph{In a bilateral manufacturer--retailer channel, how do vertical power, pricing timing, and a stage-0 regime-choice game affect CSR-driven recycling maturity, demand, and profit allocation?}

    \item \emph{In a multi-retailer EV battery recycling network, how should CSR-related reverse-logistics infrastructure be designed, and how should the resulting cost be allocated across stakeholders in a fair and economically stable way?}
\end{itemize}

The first study answers the second question by modeling CSR-driven recycling maturity in a bilateral channel. The second study answers the third question by treating EV battery recycling as a shared reverse-logistics network whose cost must be allocated among participating firms.

\section{Research Design and Case Context}
\label{sec:intro_design}

The dissertation uses EV battery recycling as a case study because it captures several features that make CSR implementation analytically difficult. First, battery recycling has a clear environmental and social rationale: unmanaged batteries create waste and pollution risks, while recycling supports circular use of critical materials \citep{Harper2019, tembo2024lithium, rezaei2025review}. Second, the activity is supply-chain intensive. Responsible recycling requires coordination between manufacturers, retailers, collection points, consolidation centers, and recycling plants. Third, the responsibility and the operational burden are not perfectly aligned. Manufacturers may face regulatory and brand-level CSR pressure, while retailers often interact with consumers and may operate the take-back interface.

Chapter~\ref{chap:literature_review} organizes the literature around two study focuses and four model blocks. Study~1 contains an operational bilateral channel game and a stage-0 meta-game. Study~2 contains a reverse-logistics network optimization model and a cooperative cost-sharing model. This structure allows the dissertation to move from CSR as a vertical channel problem to CSR as a shared network responsibility.

Chapter~\ref{chap:model1} develops Study~1. The operational model compares four channel regimes: Vertical Nash, Manufacturer Stackelberg, Retailer Stackelberg, and integration. In each regime, demand increases with CSR-driven recycling maturity, while the retailer bears a convex investment cost. The chapter then uses the operating-regime payoffs in a stage-0 meta-game to examine whether the firms coordinate or compete for leadership before operational decisions are made. This model block is central because it treats CSR coordination as an equilibrium object rather than assuming cooperation from the beginning.

Chapter~\ref{chap:model2} develops Study~2. It extends the CSR responsibility problem from a bilateral channel to a multi-retailer reverse-logistics network. The first model block formulates the physical EV battery recycling system as a mixed-integer linear programming problem. The second model block uses cooperative game theory and the Shapley value to allocate the optimized network cost. This design links operational feasibility with fairness and participation incentives.

\section{Contributions}
\label{sec:intro_contributions}

This dissertation makes four contributions.

First, it operationalizes CSR as recycling maturity in an EV battery supply chain. Rather than treating CSR as a fixed preference or a general reputation variable, the dissertation links CSR to concrete take-back and recycling capabilities: collection, sorting, traceability, recovery, and consumer trust. This framing connects CSR theory to operational supply-chain decisions.

Second, it analyzes CSR coordination under vertical power and leadership timing. The bilateral model in Chapter~\ref{chap:model1} shows how CSR-driven recycling maturity, price, demand, and profit depend on whether the channel operates under simultaneous play, manufacturer leadership, retailer leadership, or integration. This clarifies why CSR performance can vary even when the underlying technology and demand parameters are unchanged.

Third, it introduces a stage-0 meta-game for CSR regime choice. The stage-0 model asks whether the manufacturer and retailer coordinate or compete for leadership before operational pricing and recycling decisions occur. This contribution addresses a gap in models that impose channel structure exogenously. It shows how CSR coordination can fail or succeed as an equilibrium outcome.

Fourth, it extends CSR responsibility to multi-retailer reverse-logistics cost sharing. Chapter~\ref{chap:model2} connects a MILP-based EV battery recycling network with a cooperative cost game. By using the optimized network cost as the characteristic cost function, the model links physical efficiency, CSR responsibility, Shapley allocation, and coalition stability.

Together, these contributions position CSR as an implementable supply-chain governance problem. The EV battery recycling case demonstrates why CSR requires both strategic coordination and fair allocation of shared operational costs.

\section{Structure of the Thesis}
\label{sec:intro_structure}

Chapter~\ref{chap:literature_review} reviews the literature through the dissertation's two study focuses and four model blocks. It first reviews CSR-driven bilateral channel games and stage-0 regime-choice logic, then reviews EV battery reverse-logistics network design and cooperative cost-sharing models.

Chapter~\ref{chap:model1} presents Study~1. It develops a CSR-driven manufacturer--retailer game for EV battery recycling, derives operating outcomes under Vertical Nash, Manufacturer Stackelberg, Retailer Stackelberg, and integration, and then endogenizes regime choice through a stage-0 meta-game.

Chapter~\ref{chap:model2} presents Study~2. It extends the CSR responsibility problem to a multi-retailer EV battery recycling network. The chapter formulates the reverse channel as a facility-location and flow-optimization problem, defines the induced cooperative cost game, derives Shapley-based cost shares, and studies fairness, cost recovery, and coalition stability.

Chapter~\ref{chap:conclusion} synthesizes the findings of the two studies. It discusses how CSR commitments can be supported by coordination mechanisms, how vertical power and regime choice affect CSR implementation, and how shared reverse-logistics costs can be allocated when CSR responsibility becomes a multi-stakeholder network problem.
